\documentclass{wilson-summary}

\chaptertitle{Chapter 12}
\coursename{CE 271 Dynamics}
\semester{Summer 2026}

\begin{document}


% ==========================================
% Introduction / Overview
% ==========================================
\section*{Introduction to Dynamics Summer 2026}
Remember that this is a summer course so the level of detail is up to the students interested in the program, as such it is more trouble if and when you fall behind in the course.

% ==========================================
% Section 1 (e.g., 1.1)
% ==========================================
\section*{12.1: Introduction}

\subsection*{Mechanics}
% Definition sentence, then the two branches
\definition{Mechanics} studies the state of rest or motion of bodies subjected to forces.
\vspace{4pt}
\begin{tabular}{@{}p{3.2cm} p{9cm}@{}}
\toprule
\textbf{Branch} & \textbf{Description} \\
\midrule
\textbf{Statics}   & Equilibrium of a body at rest or moving at constant velocity $(\Sigma F = 0)$ \\[4pt]
\textbf{Dynamics}  & Deals with acceleration $(\Sigma F = ma)$ \\[2pt]
\quad Kinematics   & Study of the \emph{motion} of the body \\[2pt]
\quad Kinetics     & Study of the \emph{forces} causing the motion \\
\bottomrule
\end{tabular}
 


\subsection*{Problem Solving Approach}
\begin{enumerate}[leftmargin=1.6em, itemsep=2pt]
  \item Read the problem carefully.
  \item Draw any necessary diagram (\emph{free-body diagram}).
  \item Establish a coordinate system and apply the relevant principles
        (e.g.\ $\Sigma F = ma, \ldots$).
  \item Solve the equation algebraically, then apply a consistent set of units,
        and finally complete the solution numerically.
  \item Study the answer using technical judgement.
\end{enumerate}

\subsection*{A Little Bit of Math: Integration}
% The integration example with two methods
\textbf{Example.} Let $v = \dfrac{dx}{dt}$. Calculate $x(t)$ if $v(t) = 2t$ and $x(1) = 3$.
 
\medskip
\begin{minipage}[t]{0.48\textwidth}
\subsubsection*{Method 1 — Indefinite integral}
\[
  \frac{dx}{dt} = 2t \;\Rightarrow\; x(t) = t^2 + C
\]
Apply initial condition $x(1)=3$:
\[
  3 = 1^2 + C \;\Rightarrow\; C = 2
\]
\begin{keybox}
\[x(t) = t^2 + 2\]
\end{keybox}
\end{minipage}
\hfill
\begin{minipage}[t]{0.48\textwidth}
\subsubsection*{Method 2 — Definite integral}
\[
  dx = 2t\,dt
\]
\[
  \int_{3}^{x} dx = \int_{1}^{t} 2t'\,dt'
\]
\[
  x - 3 = \bigl[t'^2\bigr]_{1}^{t} = t^2 - 1
\]
\begin{keybox}
\[x(t) = t^2 + 2\]
\end{keybox}
\end{minipage}

\section*{12.2: Rectilinear Motion}
% Content goes here.

\subsection*{Subsection Title (if needed)}
Use \verb|itemize| with \verb|[nosep]| for compact lists:
\begin{itemize}[nosep]
    \item First item
    \item Second item
\end{itemize}

Use \verb|enumerate| similarly:
\begin{enumerate}[nosep]
    \item First step
    \item Second step
\end{enumerate}

\subsection*{Important Formulas}
\[
\text{Formula here}
\]

\subsection*{Example}
\textbf{Example:} Description.

\section*{12.3: Erratic Rectilinear Motion}
\section*{12.4 : General Curvilinear Motion}
\section*{12.5: Curvilinear Motion x,y,z Motion}
\section*{12.6: Motion of Projectile} 
\section*{12.7: Curvilinear n,t,b Motion}
\section*{12.8: Curvilinear r,$\theta$,z Motion}
\section*{12.9: Dependent Motion}
\end{document}
